\begin{papercard}{deng2011}
\cf{Stages} \stagetag{H} (threat and precondition trees as a model for a harm tree), \stagetag{S} (privacy properties, hard versus soft privacy), \stagetag{C} (a catalogue that can seed deductive codes). The LINDDUN privacy threat-modelling framework (Requirements Engineering 16:3--32, 2011; 30 pages). It is an engineering method for software systems, not a study of users; its worth for the thesis is the \emph{structure} (data-flow model $\to$ threat types per element $\to$ threat trees with preconditions $\to$ misuse cases $\to$ requirements $\to$ mitigations) and the explicit soft-privacy threats.

\cf{Problem and contribution} Security has systematic threat-modelling support (KAOS anti-goals, Microsoft's STRIDE with threat-tree catalogues) but privacy does not; ``STRIDE does not cover privacy threats'' \pc{deng2011}{2}{STRIDE does not cover privacy threats}. Three contributions: a methodology that maps privacy-specific threat types to the elements of a system model, ``analogous to STRIDE, an information flow--oriented model of the system is leveraged to guide the analysis and to provide broad coverage'' \pc{deng2011}{1}{Analogous to STRIDE, an information flow–oriented model of the system is leveraged to guide the analysis and to provide broad coverage}; a catalogue of privacy threat-tree patterns; and a mapping from privacy requirements to privacy-enhancing technologies (PETs).

\cf{Background: STRIDE in nine steps} (1) define use scenarios, (2) list external dependencies, (3) security assumptions, (4) external security notes, (5) data flow diagrams (DFDs) with data flows, data stores, processes, external entities and trust boundaries, (6) threat types, (7) identify threats per DFD element with threat-tree patterns, (8) risk, (9) mitigation. STRIDE is the acronym Spoofing, Tampering, Repudiation, Information disclosure, Denial of service, Elevation of privilege \pc{deng2011}{3}{STRIDE is an acronym for Spoofing, Tampering, Repudiation, Information disclosure, Denial of service, and Elevation of privilege}. Alternatives: attack trees, misuse cases (attack trees were found more effective in an experiment), KAOS anti-goals.

\cf{Privacy properties} The authors do not propose a new definition of privacy but select properties following Pfitzmann et al.'s terminology. \emph{Hard privacy} aims at data minimisation, assumes data should not be divulged to third parties, and models adversaries (service provider, data holder, outsiders) like a security system \pc{deng2011}{5}{The data protection goal of hard privacy refers to data minimization, based on the assumption that personal data are not divulged to third parties}. \emph{Soft privacy} assumes the data subject has lost control and must trust the controller's honesty and competence, so the aim is security of processing, purpose limitation and consent via policies, access control and audit \pc{deng2011}{5}{Soft privacy, on the contrary, is based on the assumption that data subject lost control of personal data and has to trust the honesty and competence of data controllers}. Properties: hard = unlinkability, anonymity and pseudonymity, plausible deniability, undetectability and unobservability, confidentiality; soft = \emph{content awareness} (users know what they disclose, data minimal and up to date, including consequences of incomplete or incorrect data, e.g.\ a patient's diabetes missing from an e-health record) and \emph{policy and consent compliance} (the system implements and enforces the stated policy and the user's consents).

\cf{LINDDUN threat types} Each letter negates a property: \textbf{L}inkability, \textbf{I}dentifiability, \textbf{N}on-repudiation (a threat for privacy, opposite of plausible deniability), \textbf{D}etectability, \textbf{D}isclosure of information, \textbf{U}nawareness (content unawareness), \textbf{N}on-compliance (policy and consent) \pc{deng2011}{8}{As shown in Table 2, the methodology considers seven types of threats}. Two soft-privacy threats: ``Content Unawareness indicates that a user is unaware of the information disclosed to the system'' \pc{deng2011}{9}{Content Unawareness indicates that a user is unaware of the information disclosed to the system}; ``Policy and consent Non-compliance means that even though the system shows its privacy policies to its users, there is no guarantee that the system actually complies with the advertised policies'' \pc{deng2011}{9}{Policy and consent Non-compliance means that even though the system shows its privacy policies to its users, there is no guarantee that the system actually complies with the advertised policies}.

\cf{Procedure} (a) Build a DFD from the use-case scenarios and mark trust boundaries; the DFD must be correct or the analysis is wrong \pc{deng2011}{9}{First, a data flow diagram (DFD) is created in correspondence with the application’s use case scenarios}. (b) Map threat types to DFD element types with Table 4: linkability, identifiability, disclosure apply to entities, flows, stores and processes (linkability and identifiability as \emph{pairs}, e.g.\ a message linkable to a sender); non-repudiation and detectability to flows, stores and processes; content unawareness to the entity; policy/consent non-compliance to the system as a whole \pc{deng2011}{10}{Policy and consent non-compliance is a threat that affects system as a whole}. Running example: a fictitious ``Social Network 2.0'' with 10 numbered threats after discarding irrelevant ones. (c) Detail each marked intersection with a \textbf{threat-tree pattern} giving the preconditions (vulnerabilities) for that threat \pc{deng2011}{11}{For each marked intersection in Table 4, a threat tree pattern exists showing the detailed preconditions for this specific threat category to materialize}; nodes are OR unless marked AND; many branches end in an information-disclosure security threat, e.g.\ linkability of a data store requires insufficient access control \emph{or} insufficient anonymisation and strong data mining. Catalogue includes linkability, identifiability (login by e-ID, password, token, biometrics), non-repudiation (deniable encryption, MACs, OTR), detectability (covert channels, steganography, dummy traffic), content unawareness (user supplies more PII than needed, or outdated data) and policy/consent non-compliance (policies tampered or wrongly updated by an administrator). (d) Document threats as \textbf{misuse cases}: a misuse case is ``a use case from the misactor's point of view'' \pc{deng2011}{18}{In particular, a misuse case can be considered as a use case from the misactor’s point of view}, with summary, assets/stakeholders/threats, primary misactor (accidental or intentional), basic and alternative flows, trigger, preconditions (the leaves of the tree). (e) Risk: likelihood times impact, with any existing technique (OWASP, DREAD, NIST 800-30, OCTAVE); privacy-specific risk assessment is called an open research challenge \pc{deng2011}{19}{The LINDDUN framework (similarly to STRIDE) is independent from the specific risk assessment technique that is used}. (f) Elicit requirements from misuse cases using Table 6, then choose mitigation: warn the user (lower risk), remove or turn off the feature (``the only way to reduce the risk to zero''), or counter with PETs \pc{deng2011}{20}{Removing or turning off the feature is the only way to reduce the risk to zero}; Table 7 maps hard and soft PETs (policy and feedback tools) to objectives.

\cf{Discussion points by the authors} Privacy threats typically concern \emph{relationships} among DFD elements, pair-wise or group-wise, whereas security concerns each component locally \pc{deng2011}{24}{Some privacy threats, in contrast to security, affect DFD elements pair-wise}; the process element matters less for privacy; hard-privacy technologies are active in research but poor in deployment, soft ones (legal compliance as driver) are state of the art with less research \pc{deng2011}{24}{Hard privacy technologies are active in research but poor in deployment}; inferred knowledge would need an extension of DFD semantics \pc{deng2011}{25}{For some privacy properties, an extension of the DFD semantics might be useful}; privacy and security objectives conflict in places (non-repudiation vs.\ plausible deniability) so both should be analysed together \pc{deng2011}{25}{it is interesting to analyze privacy threats together with security threats}.

\cf{Limits (card)} The paper validates the method only on a running example; larger case studies and an e-health system are announced as future work \pc{deng2011}{25}{As future work, we plan to apply the proposed framework to larger case studies}; the catalogue is a ``best effort'' by the authors and expected to change; it presupposes access to system design (DFD) and an adversary-centred view, so user-experienced harms such as distress, discrimination and manipulation are outside the threat types; no probability or severity scale; no reliability testing of analysts' use of the catalogue; 2011 version, later LINDDUN variants (GO, PRO) are not covered.

\cf{Use for this thesis} (1) A ready model for how to structure the harm tree of stage H: a root harm, OR/AND preconditions, evidence at the leaves (screens, policy statements, store labels), then a short narrative per path like a misuse case. (2) The soft-privacy threats map directly onto what the thesis can observe without code access: \emph{content unawareness} (what the user is told about what is disclosed) and \emph{policy and consent non-compliance} (stated policy versus observable behaviour). (3) The hard-threat vocabulary (linkability, identifiability, detectability, disclosure) can label the \emph{data-flow} side of a finding, e.g.\ advertising identifier or SDK sharing, but the thesis cannot verify them without traffic analysis and should mark them as inferred. (4) Use it as method inspiration alongside Solove's activities (what is done to data) and Citron and Solove's harm types (what results for people); do not present LINDDUN as the thesis' method.
\end{papercard}
